\chapter{結言}
\label{chap:conclusion}

\section{本論文の総括}
\label{sec:conclusion_overview}

本論文は，家庭環境における料理操作を対象として，RGB-D点群から構築した3Dモデルを物理世界の計算代理とし，対象物および作業環境の幾何学的・力学的・環境的特徴を分析した上で，その結果をロボットの動作生成および制御設計へ接続する操作設計アプローチを提示した．家庭の食材，容器，および調理器具は，形状，大きさ，材質，配置が一定ではなく，工場のような専用治具や固定環境を前提とできない．この条件下では，事前に与えた単一の軌道を再生するだけではなく，その場で取得した対象・環境モデルから，タスクの成立に必要な量を計算することが重要となる．

本論文で中心となる考え方は，センシングによって得た3Dモデルを認識処理の終点とせず，物理的分析と動作設計の入力として用いる点にある．具体的には，3D点群から表面位置，法線，断面積，体積，重心，および接触候補を抽出し，タスクに応じて軌道，目標角度，把持位置，および制御パラメータへ変換した．さらに，モデル上の分析だけでは吸収できない計測誤差や接触変動に対して，力覚または視覚フィードバックを組み合わせた．これにより，「センシング，モデリング，分析，動作計画・制御」という一連の流れを，物理的性質の異なる三つの料理操作に適用した．

提案するアプローチの適用可能性を検証するため，本論文では，皮むき動作，未知容器間の定量注ぎ，および切断中における食材の安定把持を取り上げた．皮むきでは不規則曲面への幾何学的適応と接触維持，注ぎでは容器を介した流体の間接操作，切断中把持では包丁外力と道具拘束に抗する力学的安定性を扱った．三つのタスクは対象となる物理現象が異なる一方，3Dモデルを基盤としてタスク固有の分析量を導出し，実行層へ接続するという設計原理を共有する．

\section{各研究課題で得られた成果}
\label{sec:conclusion_results}

\subsection{共通基盤としての3Dモデル取得・処理}
\label{subsec:conclusion_common_model}

第3章では，三つのタスクに共通する計測・実装基盤を構築した．PA10-7Cを用いた双腕システムとDENSO VS-050を用いた双腕システムについて，ロボット，RGB-Dカメラ，力覚センサ，グリッパ，エンドエフェクタ，およびソフトウェア構成を整理した．両実装環境では制御インタフェースが異なるが，多視点点群取得，座標変換，点群統合，ダウンサンプリング，平面除去，外れ値除去，クラスタリング，MLS平滑化，および法線推定からなる共通処理を用いた．

標準物体および3Dプリンタ製自由曲面物体による評価では，平均表面距離誤差が1~mm未満，最大寸法誤差が1.5~mm以下となった．容器の体積相対誤差は，底面が十分に取得された条件で約3\%以内であった．また，卓上撮影では不可視となる底面を，対象を把持して持ち上げた後に追加撮影する手法を示し，深底容器の体積相対誤差を$12.47\pm4.53$\%から$2.89\pm1.17$\%へ低減した．この結果は，ロボットの操作能力を観測改善へ利用することが，タスクに必要なモデル品質の確保に有効であることを示す．

以上により，後続する各タスクへ，位置・姿勢，色，法線，断面積，体積，重心，および表面接触候補を受け渡す共通の計算基盤を構築した．ただし，この評価は限定された物体と観測条件に基づくものであり，二種類のロボットシステムで実装したこと自体が，あらゆるロボットへの一般的な適用を保証するものではない．

\subsection{アプローチの検証：皮むき作業}
\label{subsec:conclusion_peeling}

第4章では，第1章で設定した第一の検証課題である皮むき動作を扱った．食材表面点群と作業台・到達可能領域の幾何関係からピーラーの基準軌跡を抽出し，表面法線の変化に基づいて軌跡を複数の実行区間へ分割した．さらに，皮むき済み領域の色と境界近傍の法線から食材の回転角を計算し，「一回皮むき動作」と「食材の回転操作」を反復する作業構成を実現した．

実行時には，点群モデルから得た軌跡を位置制御の基準とし，ピーラーと食材の接触力を力覚フィードバックで維持した．これにより，点群の局所誤差，食材表面の凹凸，および刃先の実際の切込み位置とモデル表面との差を実行時情報で補償した．りんご，きゅうり，なす，じゃがいも，にんじんを用いた13試行では，全ての対象で78\%以上の皮むき率を得た．この結果は，本実験で扱った対象と条件において，3D表面幾何から生成した動作と接触力制御の階層的な組合せが，不規則形状へ適応する皮むきに有効であることを示す．

一方，本手法は食材の支持軸を人手で与えること，色差が小さい対象では固定回転角を用いること，およびピーラーの機械的な切込み制限を利用することを前提とする．したがって，本成果は任意形状・任意材質の剥離操作全般を解決するものではなく，対象支持，表面認識，および工具機構に関する仮定の範囲で成立する．

\subsection{アプローチの検証：液体注ぐ作業}
\label{subsec:conclusion_pouring}

第5章では，直接把持できない液体を，容器という環境を介して操作する定量注ぎを扱った．第一の手法では，注ぎ先容器の3Dモデルと実行中に検出した液面高さから注いだ体積を推定し，体積誤差を用いたPD制御により注ぐ容器の角速度を調整した．第二の手法では，注ぐ容器の3Dモデルから注ぎ縁と注ぐ点を検出し，容器角度と流出体積の関係を事前計算した上で，目標体積に対応する角度へP制御で追従した．

形状特性の異なる三種類の容器を用いた実験により，注ぎ先容器を基準とする手法1では，容器3において平均符号付き誤差$-0.95$~ml，標準偏差1.16~mlを得た．注ぐ容器を基準とする手法2では，同容器において平均符号付き誤差4.96~ml，標準偏差2.57~mlとなった．手法1は実行中の液面観測を必要とする一方で高い定量性を示し，手法2は液面を観測できない配置でも実行できる一方で，表面張力，容器壁厚，および初期液量の誤差の影響を受けた．

この結果は，液体そのものの詳細な流体モデルを直接構築しなくても，液体を拘束する容器の幾何を計算代理として用いることで，注いだ体積と制御目標を導出できることを示す．ただし，手法1は液面が観測できること，手法2は注ぐ容器の内形状と初期液量が得られることを必要とし，低粘度液体，凸形状に基づく断面計算，および表面張力を無視した近似などの仮定を含む．

\subsection{アプローチの検証：切断中の安定把持の作業}
\label{subsec:conclusion_holding}

第\ref{ch:holding}章では，第1章で設定した第三の検証課題として，包丁から加わる外力に抗しながら食材を安定に保持する把持位置探索を扱った．食材表面点群と包丁軌道を入力とし，まず包丁・作業台・指間の幾何的干渉を用いて探索空間を削減した．次に，包丁刃と食材の接触から切断力およびトルクを推定し，その外力を打ち消す把持力候補を生成した．さらに，把持力の大きさ，方向，および指間分散に基づく動的評価を包丁軌道全体で累積し，最終把持位置を選定した．

ジャガイモと玉ねぎを用いた比較実験では，高い評価スコアを持つ把持位置は，ランダム候補および低スコア候補と比べて切断前後の移動量が小さい傾向を示した．また，選定した把持位置を用いた完全切断を実機で実行し，対象を保持したまま切断工程を完了できることを確認した．これにより，3D形状だけでなく，包丁軌道に沿って変化する外力と道具拘束をモデル上で評価することが，把持位置の選定に有効であることを示した．

一方，本手法は平坦な支持面と準静的な力学近似を用い，計画はオフラインで実行する．約169点の点群を用いた条件でも計算時間は75.4秒であり，切断途中の安定性評価も切断前後画像を中心としている．さらに，指先位置決め誤差と点群ノイズが外れ値を生じさせた．したがって，任意形状の食材や高速切断への拡張には，支持条件の一般化，動的センシング，および計算高速化が必要である．


\section{今後の展望}
\label{sec:conclusion_future}

今後の第一の課題は，3Dモデル取得のロバスト化と不確実性の明示である．複数センサ，偏光・熱・近接・触覚情報，および物体を動かしながら観測する能動知覚を組み合わせ，透明・反射・遮蔽領域を補完する必要がある．また，点群の各領域に位置・法線・体積推定の信頼度を付与し，信頼度が低い場合に再観測または安全側の動作へ切り替える仕組みが有効である．

第二の課題は，分析モデルのオンライン更新である．皮むきや切断では対象形状が作業中に変化し，注ぎでは液体量が連続的に変化する．実行前に固定したモデルだけでなく，作業中の視覚・力覚情報から形状，接触状態，物性値を更新し，軌道または把持計画を再計算する閉ループ化が必要である．

第三の課題は，解析モデルと学習の補完的統合である．本論文の解析は，必要な入力・出力と物理仮定が明示される利点を持つ一方，表面張力，食品の不均質性，摩擦変動などを完全には表現できない．解析モデルで大域的な構造と安全制約を与え，残差誤差，物性推定，パラメータ調整，および失敗回復をデータ駆動手法で学習する構成が考えられる．

第四の課題は，計算と実行の高速化である．特に切断中把持では，把持力候補の生成と評価が計算時間の大部分を占める．候補生成の階層化，GPU並列化，近似最適化，および過去タスクからの初期解再利用により，環境変化へ応答可能な計画速度を目指す必要がある．

第五の課題は，複数料理操作の統合と実環境評価である．同一の対象モデルを皮むき，切断，移送，注ぎへ引き継ぎ，器具の持替え，作業順序，衛生，安全，人との共有空間を含む一連の調理工程として評価する必要がある．さらに，対象種類，個体差，照明，配置，液体特性，工具の差を系統的に変化させ，成功率，精度，作業時間，失敗モード，および回復率を統計的に評価することで，提案アプローチの適用境界を明確化できる．


\chapter*{謝辞}
\addcontentsline{toc}{chapter}{謝辞}

本論文をまとめるにあたり，研究課題の設定から論文全体の構成，および研究者としての姿勢に至るまで，長期にわたり丁寧な指導をいただいた主任指導教員の工藤俊亮教授に深く感謝する．
研究の方向性に迷った際にも，多くの有意義な助言をいただき，本論文を一つの研究としてまとめることができた．

% TODO: 審査委員の氏名
本論文の審査にあたり，専門的な見地から有益な指摘と助言をいただいた審査委員の先生方に厚く御礼を述べる．
いただいた指摘は，本研究の意義と限界を明確にし，論文の内容を改善する上で大きな助けとなった．

本研究の遂行において，共同研究，実験装置の構築，ソフトウェア開発，実験の実施，および日々の議論に協力いただいた研究室の教職員，卒業生，在学生の皆様に感謝する．
多くの技術的課題を解決できたのは，皆様から得た知識，助言，および協力によるものである．

また，長期間にわたる研究生活を理解し，支え続けてくれた家族と友人に心から感謝する．
研究が思うように進まない時期にも変わらず励ましてくれたことが，本論文を完成させる大きな力となった．

最後に，本研究に関わり，直接または間接に支援いただいた全ての方々へ，改めて深い感謝を表する．

